\scanpage{561}
\mspage{285}

\[
(53)\qquad
\widetilde M_{1,1}\simeq M_{1,1}[2]'/\Gamma'_2
\simeq U_{0,3}/\unsure{\mathfrak S}_3
\]

\unsure{or}\qquad
(valable sur $\mathrm{Spec}\,\mathbb Z[1/2]$)

on trouve bien, au dessus de $\mathrm{Spec}\,\mathbb Z[1/2]$, que
$\widetilde M_{1,1}$ se déduit d'un schéma relatif $Y$, qui est
\deleted{\illegible{1}} loc. \inserted{(ét)} isom. à $\mathbb P^1_S$,
en enlevant une section $s$ (dont l'existence implique déjà que $Y$ est
\unsure{globalmt|globalement} isom. à $\mathbb P^1_S$ - donc
$Y-s(S)$ isom. à $\mathbb E^1_S$.
Dans cette approche on a donc envie de désigner par $\infty$
(\unsure{non p.} $0$, comme en théorie des cartes !) cette section,
qui correspond aux ``points à l'infini'' $\unsure{(0,1,\infty)}$ de
$U_{0,3}$, on aura un ``\unsure{pt} à l'infini'' de
$\widetilde M_{1,1}$.
Il est d'ailleurs facile de vérifier a priori \unsure{que}
(par la compactification de Deligne-Mumford de $M_{1,1}$ et de
$\widetilde M_{1,1}$) que $\widetilde M_{1,1}$ sur $\mathrm{Spec}\,\mathbb Z$
tt entier est de la forme $Y-\mathrm{Im}\,s$, où $Y$ est lisse propre
sur $\mathrm{Spec}\,\mathbb Z$, et $s$ une section - dès lors ce qu'on
connait p.ex. sur la fibre générique, et le fait que $\mathbb Z$ soit
principal, impliquent que $Y\simeq\mathbb P^1_{\mathbb Z}$, et

\scanpage{562}

qu'on a donc $\widetilde M_{1,1}\simeq\mathbb E^1_{\mathbb Z}$.
Pour \deleted{le} choisir cet isom., on utilise les deux sections de
$\mathbb E^1_{\mathbb Z}$, correspondant aux courbes anharmoniques.
Ces deux sections \uline{ne sont pas} \uline{disjointes}, elles se
rencontrent en car. $2$ et car. $3$.
(Pour ce qui se passe en car. $\ne 2$ i.e. sur
$\mathrm{Spec}\,\mathbb Z[1/2]$, on le voit bien par l'isom $53$ -
\deleted{les} car en car. $3$ les deux orbites
\[
(2,-1,1/2)\quad\text{et}\quad(j,\bar j)
\]
coïncident et se \unsure{collapsent} en un seul et
$\overline{\mathrm m}$ point).
Mais prenant l'une de ces sections comme section nulle
(sauf erreur c'est celle qui correspond à la courbe elliptique la plus
riche en symétries, donc de groupe $\mathbb Z/6\mathbb Z$, qu'on prend -
c'est une question de convention, bien sûr - c'est donc aussi l'orbite
$\{j,\bar j\}$, correspondant à un stabilisateur
$\simeq\mathbb Z/3\mathbb Z$) - cela détermine l'isom
$\widetilde M_{1,1}\simeq\mathbb E^1_{\mathbb Z}$ au signe près.
L'autre section devient alors un entier de la forme
$\pm 2^a3^b$ ($a,b\in\mathbb N^*$ - \uline{entiers} bien déterminés dont
\unsure{j'ai} \illegible{3} la valeur par \unsure{cœur}), et on achève de

\scanpage{563}
\mspage{286}

normaliser en exigeant que le \uline{signe} soit \uline{plus}. Donc

\uline{Théorème} (pour mémoire - c'est bien connu) Il y a un isom unique
\[
(54)\qquad
\widetilde M_{1,1}\simeq\mathbb E^1_{\mathbb Z},
\]
qui \deleted{donne} en car. $0$ donne à la courbe anharmonique de groupe
d'autom $\mathbb Z/6\mathbb Z$ le
\deleted{\illegible{1}}\inserted{\unsure{j-invariant}} $0$, et à celle de
groupe $\mathbb Z/4\mathbb Z$ un inv. \deleted{\illegible{1}} $>0$
(qui sera nécessairement $2^a3^b$, avec $a,b\in\mathbb N^*$ bien
déterminés, mais par moi oubliés\ldots).
\margin{$b=3$, cf ci-dessous}

En \unsure{termes} de l'isom
\[
\deleted{\unsure{(55)}}\qquad
\widetilde M_{1,1}\simeq U_{0,3}/\unsure{\mathfrak S}_3
\subset \mathbb P^1/\unsure{\mathfrak S}_3
\]
(valable sur $\mathrm{Spec}\,\mathbb Z[1/2]$)
ceci correspond à \deleted{l'isom} un isom entre
$\mathbb P^1/\unsure{\mathfrak S}_3$ et $\mathbb P^1$ qui à l'orbite
$(0,1,\infty)$ associe $\infty$ (non $0$) et à l'orbite $(j,\bar j)$
associe $0$ (non $\infty$), à l'orbite $(2,-1,1/2)$
\unsure{le fameux} $2^a3^b$ (et non $1$). La fonction
\margin{\uline{invariant} \uline{modulaire}}
\[
(55)\qquad j(\lambda)\in\mathbb Q(\lambda)
\]
qui réalise cet isom, i.e \unsure{l'homomorphisme}
\[
(56)\qquad
\mathbb P^1_{\mathbb Z}\xrightarrow{j}\mathbb P^1_{\mathbb Z}
\]
correspondant (de degré $6$, avec
\[
j\circ\unsure{u|\mu}=j
\qquad(\unsure{u|\mu}\in\unsure{\mathfrak S}_3)
\]
)

\scanpage{564}

est donc \unsure{liée} à celle inspirée de la théorie des cartes,
soit $f(\lambda)$, par
\[
(57)\qquad
j(\lambda)=2^a3^b f(\lambda)^{-1}.
\]

D'ailleurs $f(\lambda)$ \inserted{(ou $j(\lambda)$)} se calcule aisément
par connaissance de ses zéros et de ses pôles,
\inserted{avec leurs multiplicités},
\deleted{\illegible{2}} et la ``normalisation'' par la valeur de $f$
(ou de $j$) sur une des sections $2,-1,1/2$ ; on trouve immédiatement :
\[
(58)\qquad
f(z)=\frac{3^3}{2^2}\frac{z^2(z-1)^2}{(z^2-z+1)^3}
\]
donc
\[
(59)\qquad
j(\lambda)=2^{a+2}3^{b-3}\frac{(z^2-z+1)^3}{z^2(z-1)^2}.
\]

Comme $j(\lambda)$ doit garder un sens en car. $3$ et ne pas être
constante, cela montre d'ailleurs \uline{$b=3$}, donc $(59)$ s'écrit aussi
\[
\deleted{\unsure{(60)}}\qquad
j(\lambda)=2^{a+2}\frac{(z^2-z+1)^3}{z^2(z-1)^2}
\]
- mais on se rappelle que cette formule n'est pertinente -
\unsure{ne permet} de calculer
\deleted{un \illegible{1} de}\inserted{l'invariant d'une}
courbe elliptique - qu'en car $\ne 2$, heureusement !
Pour \deleted{\illegible{1}}

\scanpage{565}
\mspage{287}

\deleted{définir}\inserted{décrire}
\[
(60)\qquad
M_{1,1}\longrightarrow\mathbb E^1_{\mathbb Z}
\]
\margin{et en particulier pour déterminer
\unsure{l'entier invariant|l'autre invariant} \illegible{2}
$2^a3^3$, i.e. \uline{déterminer $a$}}
aussi au voisinage de $2$ - disons en car. $\ne 3$ - il faut une étude
des courbes elliptiques qui ne passe plus par la représentation
(de \uline{Legendre}, sauf erreur)
$y=\sqrt{x(x-1)(x-\lambda)}$, ou encore, par la rigidification des
\unsure{points|pts} d'ordre $2$, mais par exemple par celle d'ordre $3$.

\uline{Remarque} On voit tout de suite que le quotient
$\mathbb P^1_{\mathbb Z}/\unsure{\mathfrak S}_3$ s'identifie à
$\mathbb P^1_{\mathbb Z}$, de façon unique en prenant comme images des
orbites $(0,1,\infty)$ et $(j,\bar j)$ les sections $0$ et $\infty$
(p.ex., en suivant la convention qui correspond à la théorie des cartes),
et en prenant comme image de l'orbite $(2,-1,1/2)$ un
\deleted{entier}\inserted{rationnel} qui soit $>0$.
Ceci est possible a priori, car on sait que les \deleted{sections}
orbites $(0,1,\infty)$ et $(j,\bar j)$
\deleted{d'une part, $(0,1,\infty)$ et $(2,-1,1/2)$ d'autre part,}
ne coïncident en aucune caractéristique, donc correspondent à des sections
disjointes de $Y=\mathbb P^1_{\mathbb Z}/\unsure{\mathfrak S}_3$,
tandis que l'orbite \deleted{\unsure{restante}}

\scanpage{566}

$(2,-1,1/2)$ coïncide avec la première en car. $2$, avec la deuxième en
car. $3$ (et ce sont là les seules coïncidences qui peuvent arriver).
On sait donc a priori que la troisième section sera de la forme
$2^\alpha/3^\beta$, avec $\alpha,\beta\in\mathbb N^*$.
Le calcul explicite est d'ailleurs évident, l'hom composé
\[
(61)\qquad
\mathbb P^1_{\mathbb Z}
\xrightarrow{f_1^{\unsure{*|\circ}}}
\mathbb P^1_{\mathbb Z}/\unsure{\mathfrak S}_3
\simeq\mathbb P^1_{\mathbb Z}
\]
doit être de la forme
$f_1^{\unsure{*|\circ}}=c.f^{\unsure{*|\circ}}$, avec la constante
\inserted{$c\in\mathbb Q$} choisie de telle façon que
$cf^{\unsure{*|\circ}}$ \deleted{ait} se réduise bien en
\unsure{ttes} car. (ce qui détermine $c$ \illegible{2} \uline{signe})
et que
\inserted{$f_1(-1)=$}\deleted{\illegible{1}}
$cf^{\unsure{*|\circ}}(-1)$ \inserted{$=c$} $>0$
(ce qui fixe l'indétermination du \uline{signe}).
On trouve alors
\[
(62)\qquad
f_1(\lambda)=
\frac{\lambda^2(\lambda-1)^2}{(\lambda^2-\lambda+1)^{\unsure{2|3}}}
=\frac{2^2}{3^3}f(\lambda)
\]
donc
\[
(63)\qquad
f_1(\unsure{-1|2})=\frac{2^2}{3^3}
\qquad(\text{donc }\alpha=2,\ \beta=3).
\]

Mais on fera attention que ces calculs
\inserted{\unsure{faits au voisinage de la car. $2$}}
ne se rapportent plus à $\widetilde M_{1,1}$ et
$\widetilde{\widetilde M}_{1,1}$,
où la configuration des trois sections
\deleted{\illegible{1}
$\unsure{\Delta_\infty,\Delta_0|\Lambda_\infty,\Lambda_0}$,}
n'est pas la $\overline{\mathrm m}$ \illegible{1}\ldots

\scanpage{567}
\mspage{288}

Sur la \unsure{lancée} de ces réflexions, il serait naturel de regarder
de plus près le schéma
\[
(64)\qquad
M_{1,1}[4],\quad
\text{sur lequel agit }\Gamma_4=\mathrm{Gl}(2,\mathbb Z_4)
\]
\[
\text{ou}\quad \Gamma'_4=\Gamma_4/\pm1
\]
\margin{\illegible{15}}
et \deleted{\illegible{1}} \unsure{ci-dessous}
\[
(65)\qquad
M_{1,1}[4]\longrightarrow M_{1,1}[2]'\simeq U_{0,3}
\]
\unsure{(compatible avec l'hom}
\[
\Gamma'_4=\mathrm{Gl}(2,\mathbb Z/4)/\pm1
\longrightarrow
\Gamma'_2=\Gamma_2\simeq\mathrm{Gl}(2,\mathbb Z/2)
\simeq\unsure{\mathfrak S}_3
\]
\margin{\diagram{[Au-dessus du terme de gauche, un petit trait fléché
oblique vers la gauche, près de l'inscription
\deleted{\unsure{$\det$}}.]}}
\margin{\unsure{on fait le noyau de l'hom}
\illegible{35}
$\unsure{\mathbb F_2}$
\illegible{4}
$\unsure{\mathbb F_2^2}$}

Un calcul immédiat, compte tenu que
\[
(\mathbb Z/4\mathbb Z)^*\simeq\{\pm1\}
\]
et \unsure{que}
\[
\mathrm{card}\,\Gamma_4
=
2\underbrace{\mathrm{card}\,\mathrm{Sl}(2,\mathbb Z/4\mathbb Z)}
_{8\,\mathrm{card}\,\mathrm{Sl}(2,\mathbb Z/2\mathbb Z)}
=
16.\underbrace{\mathrm{card}\,\Gamma_2}_{6}
\]
donc
\[
(66)\qquad
\mathrm{card}\bigl(\Gamma'_{4,2}
=\mathrm{ker}(\Gamma'_4\longrightarrow\Gamma'_2)\bigr)=8
\]
i.e. $M_{1,1}[4]$ est un rev. galoisien de $M_{1,1}[2]'$,
de \deleted{groupe} ordre $8$.
Notons qu'il n'est pas géom. connexe sur $\mathbb Q$,
\inserted{\unsure{ou sur $\mathrm{Spec}\,\mathbb Z[1/2]$}},
- puisque $M_{1,1}[4]$ se trouve sur l'extension quadratique
$\mu_4^*$ \deleted{\unsure{$\simeq\mathrm{Spec}$}}, de
$\mathrm{Spec}\,\underbrace{\mathbb Z[1/2]}_{\Lambda}$ disons,
qui n'est autre que $\Lambda[i]$.
\deleted{\unsure{Choisissons le point} \illegible{1}}
\inserted{\deleted{\unsure{Considérons} \illegible{2}}}
\inserted{\unsure{factorisation de}
$M_{1,1}[4]\longrightarrow M_{1,1}[2]'$}
\deleted{\unsure{au dessus de}
$\mathrm{Spec}\,\Lambda[i]=\mu_4^*$
sur $\mathrm{Spec}\,\mathbb Z[1/2]$,}
nous pouvons distinguer un \deleted{\unsure{composant connexe}}

\scanpage{568}

\((67)\)
\diagram{
\begin{tikzcd}[column sep=large,row sep=large]
M_{1,1}[4]
  \arrow[r,"\text{galoisien}"]
  \arrow[d]
&
(U_{0,3})_{\unsure{S'}}
  \arrow[r,"\text{galoisien }\unsure{\text{degr. }2}"]
  \arrow[dl]
&
M_{1,1}[2]'\simeq U_{0,3}
  \arrow[d]
\\
S'=\mu_4^*
  \arrow[rr,"\text{galoisien}","\unsure{\text{degr. }2}"']
&&
\mathrm{Spec}\,\mathbb Z[1/2]=S
\end{tikzcd}
}

où maintenant
\[
M_{1,1}[4]\longrightarrow M_{1,1}[2]'_{\unsure{S'}}
=(U_{0,3})_{S'}
\]
est un rev. galoisien de groupe le noyau de
$\Gamma'_{4,2}\xrightarrow{\det}\{\pm1\}$, ou encore celui de
$\mathrm{Sl}(2,\mathbb Z/4\mathbb Z)'
\longrightarrow\mathrm{Sl}(2,\mathbb Z/2\mathbb Z)'$,
\deleted{\illegible{1}}
que nous allons désigner par $S\Gamma'_{4,2}$, qui est d'ordre $4$.

\uline{Proposition} Le groupe
\[
(68)\qquad
S\Gamma_{4,2}
=
\mathrm{Ker}\bigl(
\mathrm{Sl}(2,\mathbb Z/4\mathbb Z)
\longrightarrow
\mathrm{Sl}(2,\mathbb Z/2\mathbb Z)
\bigr)
\]
est isom. à \deleted{\illegible{2}}
$\unsure{\mathfrak{sl}}(2,\mathbb F_2)$
(groupe des matrices à trace nulle à coeff. dans $\mathbb F_2$),
\inserted{avec \unsure{l'op. évidente|les op. évidentes}
de $\mathrm{Sl}(2,\mathbb F_2)$ dessus,}
et le sous-groupe $\pm1$ correspond aux matrices scalaires -
le quotient \inserted{$S\Gamma'_{4,2}$ des deux}
\unsure{est} isomorphe à \deleted{\illegible{1}} $\mathbb F_2^2$
(de façon \unsure{idiote} et pas canonique du tout !).
\margin{\uline{NB} \unsure{(les $4$} \illegible{7}
$\unsure{\mathbb F_2\subset\mathfrak{gl}(2,\mathbb F_2)|
\mathbb F_2\subset\mathfrak{sl}(2,\mathbb F_2)}$
\unsure{(matrices scalaires)}
\illegible{30}
$\unsure{\mathbb P^1(\mathbb F_2)}$
\illegible{7}}
\deleted{\unsure{enfin}}
Regardant $M_{1,1}[4]$ comme un rev. galoisien de
$M_{1,1}[2]'_{S'}=(U_{0,3})_{S'}$, on trouve
\uline{géométriquement} le ``revêtement \uline{abélien} universel de
$U_{0,3}$ de groupe annulé par $2$ (et pour cause), qui a comme groupe
de Galois $(\mathbb F_2)^{\{0,1,\infty\}}/\mathrm{Diagonale}$.

\scanpage{569}
\mspage{289}

Enfin, tout est essentiellement tautologique - mais un calcul amusant à
faire sera de ``décrire l'action de $\mathrm{Sl}(2,\mathbb F_2)$ sur
$\unsure{\mathfrak{sl}}(2,\mathbb F_2)/\text{matrices scalaires}$''
de façon à l'identifier à l'action de $\unsure{\mathfrak S}_3$ sur
$\mathbb F_2^{\{0,1,\infty\}}/\text{diagonale}$.

Comme revêtement de \deleted{\illegible{3}}
\[
(69)\qquad
\widetilde M_{1,1\,S'}
\simeq M_{1,1}[4]/S\Gamma'_4
\simeq M_{1,1}[2]'_{\unsure{S'}}/\Gamma'_2,
\]
$M_{1,1}[4]$ est un rev. galoisien d'ordre $24$ ($=4.6$),
de groupe $\mathrm{Sl}(2,\mathbb Z/4\mathbb Z)'$.
\deleted{C'est} Ce serait bien venu \unsure{qu'on n'ait pas} un isom
\[
(70)\qquad
\mathrm{Sl}(2,\mathbb Z/4\mathbb Z)'
\simeq\unsure{\mathfrak S}_4
\qquad
(\text{automorphismes de l'octaèdre})
\]
et que $M_{1,1}[4]_{S'}$ en soit le revêtement \unsure{étale}
\deleted{\illegible{1}}
de $\mathbb P^1_{S'}\setminus\{0,1,\infty\}$, avec
\unsure{ramification|ramifications} compatible avec
$2\{1\}$ \deleted{\illegible{1}} $+3\{\unsure{\infty}\}$,
qui correspond à l'\uline{octaèdre} \unsure{\&} dans la théorie des
cartes - on devrait avoir sur
$S'=\mathrm{Spec}\,\mathbb Z[1/2][i]$
un isom. \inserted{canonique}
\deleted{\unsure{de $\widetilde M_{1,1}$}}
\[
(71)\qquad
M_{1,1}[4]_{S'}\simeq
\text{courbe de Fermat}
\quad
x^2+y^2+z^2=0
\quad(\text{sur }S')
\]
\margin{\unsure{il vaudrait mieux} écrire l'équation de Fermat
\unsure{ainsi}
$x_0^2+x_1^2+x_\infty^2=0$}
(dont le groupe d'autom. est \unsure{justement} une extension de
$\unsure{\mathfrak S}_3$ par
$\mu_2^{\{0,1,\infty\}}/\mathrm{diag}$ !).

\scanpage{570}

Pour bien faire, il faudrait expliciter la relation entre courbes
elliptiques \inserted{$E/S$} (sur un schéma $S$ au dessus de
$\mathbb Z[1/2][i]$) \unsure{munies} d'une rigidification de Jacobi
d'échelon $4$ $e_1,e_2$ avec $e_1\wedge e_2=[i]^{-1}$, et solutions
``non triviales'' de l'équation
\[
(72)\qquad x^2+y^2+z^2=0
\]
sur $S$ (i.e les systèmes de sections
$x,y,z\in\Gamma(S,\mathcal O_S^*)$ satisfaisant $(72)$, modulo
multiplication par un ``scalaire''
\inserted{$\alpha\in\Gamma(S,\mathcal O_S^*)$}
en commun, \unsure{rompant} la symétrie en posant
$X=x/z$, $Y=y/z$, \deleted{\illegible{1}}
les solutions de l'équation
\[
(73)\qquad X^2+Y^2=-1\qquad(=i^2)
\]
$(X,Y\in\Gamma(S,\mathcal O_S^*))$\ldots

On a ainsi interprété $M_{1,1}[2]'$, et
(modulo des vérifications et une étude un peu plus poussée)
$M_{1,1}[4]$, en relation avec \deleted{deux} cartes
\inserted{\unsure{triangulaires} \illegible{1}}
\unsure{régulières} - la carte ``\unsure{dihédrale|dièdrale}''
\inserted{\unsure{d'indice $3$|d'ordre $3$}}
ou carte triangulaire \illegible{1} universelle
(trois sommets $0,1,\infty$,
\figure{[Un contour presque circulaire, avec un ovale horizontal.
Trois points pleins sont marqués sur l'ovale : $0$ à gauche,
$\infty$ sur la branche supérieure à droite, $1$ sur la branche
inférieure à droite. Les trois portions de l'ovale joignent ces
trois points.]}
trois arêtes, deux faces qui sont des triangles, type
$(p,q)=(3,2)$), et la carte \uline{\unsure{octogonale}},
qui est un rev. d'ordre $4$ de celle-ci
\unsure{(}\deleted{\illegible{2}}\inserted{\unsure{étale}}
en dehors des sommets $\{0,1,\infty\}$),

Il y aurait \unsure{lieu} aussi de regarder des
